\documentclass[10pt,a4paper]{amsart}
\usepackage[margin=19mm]{geometry}
\usepackage{amsmath,amssymb,mathtools}
\usepackage[scr=rsfs,cal=euler]{mathalfa}
\usepackage{xfrac}
\usepackage[unicode,hidelinks,bookmarks=false]{hyperref}
\newcommand{\ie}{i.e.\ }
\newcommand{\cf}{cf.\ }
\newcommand{\resp}{respectively\ }
\newcommand{\Subsection}{\subsection}
\providecommand{\nobreakdash}{\nobreak\hbox{-}}
\setlength{\emergencystretch}{2em}
\multlinegap0pt
\usepackage{bm}
\usepackage{bbm}
\usepackage{dsfont}
\usepackage{xspace}
\usepackage{cancel}
\usepackage{mathtools}
\usepackage{cleveref}
\usepackage{amsthm}
\makeatletter
\@ifundefined{lemma}{\newtheorem{lemma}{Lemma}}{}
\@ifundefined{remark}{\newtheorem*{remark}{Remark}}{}
\makeatother
\newcommand{\N}{\mathbb{N}}
\newcommand{\abs}[1]{\left\lvert#1\right\rvert}
\newcommand{\norm}[1]{\lVert#1\rVert}
\newcommand{\schedule}{\bm{\eta}}
\title[Lower Bounds for Anytime Acceleration of Gradient Descent]{Lower Bounds for Anytime Acceleration of Gradient Descent}
\author{Chung-En Tsai, Ilyas Fatkhullin, Liang Zhang, Niao He}
\date{Source: arXiv:2607.02053v1 (CC BY 4.0)}
\begin{document}
\maketitle

\noindent\emph{Source and licence.} Lower Bounds for Anytime Acceleration of Gradient Descent. Version arXiv:2607.02053v1 (CC BY 4.0), distributed under the Creative Commons Attribution 4.0 International licence, as recorded in the arXiv licence metadata for this exact version and shown on the abstract page https://arxiv.org/abs/2607.02053v1. Full rights notice: licenses/arxiv-2607.02053-CC-BY-4.0.txt.\par\smallskip

\noindent\textit{Editorial scope.} Counted unit: the Remark whose source label is \texttt{None} in the article (source0.tex lines 847--853, source label \texttt{None}), delivered together with the article's own complete original proof of it (source0.tex lines 855--913, one \texttt{proof} environment of 59 lines). No mathematics is written, repaired or re-derived: statement and proof are the article's own text line for line. Required context retained uncounted: eq:quadratic\_lower\_bound (lines 469--472); eq:asymmetric\_huber (lines 735--741); lem:asymmetric\_huber (lines 826--841); eq:asymmetric\_huber\_part\_i (lines 828--833); eq:asymmetric\_huber\_part\_ii (lines 835--840). Every definition, display and label of those lines is retained unchanged. Omitted from this package and not redistributed: 1--468, 473--734, 742--825, 834--834, 841--846, 854--854, 914--1303. Nothing inside a retained excerpt is omitted or altered: every retained body is delivered exactly as the article's lines are written, so no omission receipt of any kind accompanies this package. Citations: the retained excerpts carry no citation command at all, so no bibliography entry of the article is transcribed and no reference is redistributed. Reference adaptation: none was needed and none was made; every reference inside the delivered unit resolves inside it, so no unresolved reference or citation is produced. Printed slips kept literal: no printed word, symbol, hypothesis or number of the counted statement or of its proof was altered. Layout and class: the article's own class is \texttt{article} with packages including inputenc, fontenc, authblk, hyperref, url, booktabs, graphicx, amssymb, amsmath, amsthm, mathtools, mathrsfs, braket, bbm; the wrapper is the lane's pinned \texttt{amsart} editorial wrapper at 10pt on A4, which restates the theorem-like environments the retained text needs and adds the article's own macro definitions that the retained text needs (\texttt{\textbackslash N}, \texttt{\textbackslash abs}, \texttt{\textbackslash norm}, \texttt{\textbackslash schedule}), with extra wrapper packages (bm, bbm, dsfont, xspace, cancel, mathtools, cleveref). The delivered numbering is the unit's own. Assets: no retained span contains a figure environment, a graphics-inclusion command, a PDF-page inclusion, a table or a TikZ picture; no image file is copied into this package, the package declares no asset dependency and no image file is redistributed with it. Licence: version arXiv:2607.02053v1 of this article is distributed under the Creative Commons Attribution 4.0 International licence, as recorded in the arXiv licence metadata for this exact version and shown on the abstract page https://arxiv.org/abs/2607.02053v1. A full rights notice is delivered at licenses/arxiv-2607.02053-CC-BY-4.0.txt. Characters: every retained excerpt is pure ASCII, so the delivered engine, XeLaTeX with its default Latin Modern text font, reports no missing character.
\begin{equation}\label{eq:quadratic_lower_bound}
	Q_n(\schedule)
	\coloneqq \max_{\lambda\in[0,1]}\lambda \prod_{k=1}^n (1 - \eta_k \lambda )^2.
\end{equation}

\begin{equation}\label{eq:asymmetric_huber}
	f(x) \coloneqq \begin{cases}
		\delta\abs{x} - \frac{ \delta^2 }{2} &\text{if }x\geq\delta, \\
		\frac{1}{2}x^2 &\text{if }-\varepsilon \leq x \leq \delta, \\
		\varepsilon\abs{x} - \frac{ \varepsilon^2 }{2} &\text{if }x\leq-\varepsilon.
	\end{cases}
\end{equation}

\begin{lemma}\label{lem:asymmetric_huber}
For any $\schedule\in(0,\infty)^{\N}$, any $n\in\N$, and any $m\in[n]$ such that $\eta_m > 1$, we have
\begin{equation}\label{eq:asymmetric_huber_part_i}
\begin{split}
	R_n(\schedule) &\geq \frac{ (\eta_m-1)^2 }{ ( 1 + \eta_{1:m-1} )^2 ( 1 + 2 \eta_{m+1:n} ) }, \\
	G_n(\schedule) &\geq \frac{ (\eta_m-1)^2 }{ ( 1 + 2 \eta_{1:m-1} ) ( 1 + \eta_{m+1:n} )^2 }.
\end{split}
\end{equation}
In particular, we have for any $n\in\N$,
\begin{equation}\label{eq:asymmetric_huber_part_ii}
\begin{split}
	\eta_n &\leq 1 + (1 + \eta_{1:n-1})\sqrt{R_n(\schedule)}, \\
	\eta_n &\leq 1 + \sqrt{1 + 2\eta_{1:n-1}}\sqrt{G_n(\schedule)}.
\end{split}
\end{equation}
\end{lemma}

\begin{remark}
\Cref{lem:asymmetric_huber} possesses two interesting properties.
The lower bounds for $R_n$ and $G_n$ differ, and they are \emph{not} permutation-invariant:
Permuting the first $n$ stepsizes may change the values of the lower bounds.
This stands in contrast to the lower bounds derived from the convex quadratic functions \eqref{eq:quadratic_lower_bound}.
There, the two lower bounds are identical and permutation-invariant.
\end{remark}

\begin{proof}
We consider the asymmetric Huber function \eqref{eq:asymmetric_huber} with $\delta= \frac{ 1 }{ 1+\eta_{1:m-1} } > 0$ and some
\begin{equation}\label{eq:eps_range}
	0< \varepsilon \leq \frac{  (\eta_m-1)\delta }{ 1 + \eta_{m+1:n} },
\end{equation}
which will be determined later.
Since $\eta_m>1$ and $\delta>0$, the range of $\varepsilon$ is nonempty.
It can be checked that $f$ is $1$-smooth convex and has a unique minimizer at $x^\star=0$.

Now we let $x_1 = 1 > \delta$ be the initial point of GD and compute its trajectory.
By induction, it can be shown that for all $k\in[m]$, we have $x_k\geq\delta$ and $x_k = 1 - \eta_{1:k-1}\delta$.
In particular, $x_m = 1 - \eta_{1:m-1}\delta = \delta$ by the definition of $\delta$.
In the $m$-th step, $x_{m+1} = x_m - \eta_m f'(x_m) = \delta - \eta_m\delta = -(\eta_m-1)\delta$.
Note that
\begin{align*}
	x_{m+1} = -(\eta_m-1)\delta \leq -\frac{  (\eta_m-1)\delta }{ 1 + \eta_{m+1:n} } \leq -\varepsilon,
\end{align*}
where the last step follows from \eqref{eq:eps_range}.
Lastly, by induction, we have $x_k\leq-\varepsilon$ and $x_{k} = x_{m+1} + \eta_{m+1:k-1}\varepsilon$ for every $k=m+1,\ldots,n+1$.


We are now ready to derive lower bounds on $R_n$ and $G_n$.
For $R_n$, by completing the square with respect to $\varepsilon$, we have
\begin{align*}
	f(x_{n+1})
	&= \varepsilon\abs{ x_{n+1} } - \frac{ \varepsilon^2 }{2}  \\
	&= -\varepsilon( x_{m+1} + \eta_{m+1:n}\varepsilon ) - \frac{ \varepsilon^2 }{2} \\
	&= - \frac{1}{2} \left( 1 + 2\eta_{m+1:n} \right) \varepsilon^2 - \varepsilon x_{m+1} \\
	&= - \frac{1}{2} \left( 1 + 2\eta_{m+1:n} \right) \left( \varepsilon + \frac{x_{m+1}}{1+2\eta_{m+1:n}} \right)^2 + \frac{ x_{m+1}^2 }{ 2 ( 1+2\eta_{m+1:n} ) }.
\end{align*}
By choosing $\varepsilon = -\frac{ x_{m+1} }{ 1 + 2\eta_{m+1:n} } = \frac{  (\eta_m-1)\delta }{ 1 + 2\eta_{m+1:n} } \leq \frac{  (\eta_m-1)\delta }{ 1 + \eta_{m+1:n} }$ that maximizes $f(x_{n+1})$, we obtain
\[
	f(x_{n+1}) = \frac{ x_{m+1}^2 }{ 2 ( 1+2\eta_{m+1:n} ) }
	= \frac{ (\eta_m-1)^2 }{ 2 (1 + \eta_{1:m-1})^2( 1+2\eta_{m+1:n} ) }.
\]
Since $\abs{x_1-x^\star}=1$, we conclude that
\[
	R_n(\schedule) \geq \frac{ (\eta_m-1)^2 }{ (1 + \eta_{1:m-1})^2( 1+2\eta_{m+1:n} ) }.
\]
For $G_n$, we have
\[
	G_n(\schedule)
	\geq \frac{ \frac{1}{2}\norm{ \nabla f(x_{n+1}) }^2 }{ f(x_1) - f(x^\star) }
	= \frac{ \varepsilon^2 }{ 2\delta - \delta^2 }.
\]
By choosing $\varepsilon = \frac{  (\eta_m-1)\delta }{ 1 + \eta_{m+1:n} }$ to maximize the lower bound, we obtain
\[
	G_n(\schedule)
	\geq \frac{ (\eta_m-1)^2\delta^2 }{ (2\delta-\delta^2)(1+\eta_{m+1:n})^2 }
	= \frac{ (\eta_m-1)^2 }{ (1 + 2\eta_{1:m-1})(1+\eta_{m+1:n})^2 },
\]
where the last equality follows from the definition of $\delta=\frac{1}{1+\eta_{1:m-1}}$.
This proves \eqref{eq:asymmetric_huber_part_i}.

It remains to prove \eqref{eq:asymmetric_huber_part_ii}.
If $\eta_n\geq 1$, then \eqref{eq:asymmetric_huber_part_ii} follows by taking $m=n$ in \eqref{eq:asymmetric_huber_part_i} and rearranging the inequalities.
If $\eta_n\leq 1$, then \eqref{eq:asymmetric_huber_part_ii} holds trivially.
This completes the proof.
\end{proof}

\end{document}
